\chapter*{Giới thiệu: AI tạo sinh trong học thuật - Một biên giới mang tính chuyển đổi}
\addcontentsline{toc}{chapter}{Giới thiệu: AI tạo sinh trong học thuật - Một biên giới mang tính chuyển đổi}
\markboth{Giới thiệu: AI tạo sinh trong học thuật - Một biên giới mang tính chuyển đổi}{Giới thiệu: AI tạo sinh trong học thuật - Một biên giới mang tính chuyển đổi}

\subsection{Tóm tắt}

Phần giới thiệu đặt nền tảng cho một cuộc khảo sát mang tính phê phán và thực tiễn về vai trò của AI tạo sinh trong môi trường học thuật. Phần này lập luận rằng các công cụ như GPT, Claude, Gemini và Llama không chỉ đơn thuần tự động hóa các tác vụ học thuật mà còn làm biến đổi bản chất của nghiên cứu, viết lách, giảng dạy và thực tiễn thể chế. Thông qua việc phác họa quỹ đạo lịch sử của AI, từ logic biểu tượng đến các mô hình ngôn ngữ lớn, phần này nhấn mạnh ý nghĩa của bước ngoặt tạo sinh và những hệ quả của nó đối với việc sản sinh tri thức. Chương này nêu các mục tiêu của cuốn sách: cung cấp cho các nhà nghiên cứu, nhà giáo dục và quản trị viên học thuật một khung tham chiếu cân bằng, dễ tiếp cận và dựa trên nghiên cứu để tiếp cận AI một cách có trách nhiệm. Các khái niệm nền tảng, kiến trúc mô hình và cơ chế sinh ngôn ngữ được giới thiệu nhằm trang bị cho độc giả sự hiểu biết về mặt kỹ thuật và lịch sử đối với những công cụ đang định hình lại công việc học thuật.

\subsection{Từ khóa}

AI tạo sinh, nghiên cứu học thuật, mô hình ngôn ngữ lớn, đạo đức AI, giáo dục đại học

Trí tuệ nhân tạo không còn là một khái niệm lý thuyết xa vời, nó đang tích cực định hình lại công việc học thuật. Sự xuất hiện của AI tạo sinh, đặc biệt là các mô hình ngôn ngữ lớn (LLM) như GPT của OpenAI, Claude của Anthropic, Llama của Meta và Gemini của Google, đã mở ra những cách thức mới để tiến hành nghiên cứu, tạo nội dung, phân tích dữ liệu và thậm chí tương tác với sinh viên. Khi các công nghệ này tiếp tục phát triển, chúng không chỉ là công cụ tự động hóa mà còn là chất xúc tác cho sự chuyển đổi, tác động đến tư duy học thuật, thực tiễn thể chế và các cuộc tranh luận về đạo đức.

Cuốn sách này khảo sát tiềm năng chuyển đổi của AI tạo sinh trong môi trường học thuật, đưa ra những hiểu biết thực tiễn về các ứng dụng, cơ chế vận hành và các hàm ý rộng hơn của nó. Bằng cách phân tích việc triển khai AI hiện nay và dự báo những diễn biến trong tương lai, chúng tôi trang bị cho học giả, nhà giáo dục và nhà quản lý những kiến thức thiết yếu để định hướng, và phát triển mạnh trong, bối cảnh đang thay đổi nhanh chóng này.

\section{Mục đích và phạm vi}

Mục tiêu chính của cuốn sách này là cung cấp một góc nhìn có cấu trúc, dựa trên nghiên cứu về vai trò của AI tạo sinh (generative AI) trong môi trường học thuật. Chúng tôi vượt ra ngoài những lối diễn giải đơn giản hóa, vốn coi AI hoặc là giải pháp không tưởng, hoặc là mối đe dọa hiện sinh. Thay vào đó, chúng tôi khám phá những thực tế nhiều sắc thái của việc tích hợp AI vào nghiên cứu, viết lách, quản lý dữ liệu, giảng dạy và quản trị học thuật.

Để đạt được điều này, chúng tôi trình bày một lộ trình toàn diện bao gồm:

\begin{itemize}
\item
  \textbf{Sự tiến hóa của AI}: Tổng quan lịch sử về trí tuệ nhân tạo, từ các mô hình biểu tượng ban đầu đến cuộc cách mạng học sâu đã tạo nền tảng cho AI tạo sinh ngày nay (Chương 1).
\item
  \textbf{AI cho nghiên cứu}: Cách AI tạo sinh hỗ trợ tổng quan tài liệu, hình thành giả thuyết và phát triển khái niệm, đồng thời làm dấy lên những lo ngại về thiên kiến và tính liêm chính trí tuệ (Chương 2).
\item
  \textbf{Viết và công bố}: Vai trò của AI trong việc soạn thảo bài báo nghiên cứu, cấu trúc lập luận và hợp lý hóa quy trình xuất bản học thuật (Chương 3).
\item
  \textbf{Quản lý và phân tích dữ liệu}: Việc tích hợp AI vào cấu trúc hóa, làm sạch và diễn giải các tập dữ liệu lớn, giảm khối lượng công việc thủ công và mở ra những hình thức khám phá mới (Chương 4).
\item
  \textbf{Trình bày và phổ biến}: Cách AI nâng cao chất lượng các bài thuyết trình học thuật, sự tương tác với công chúng và việc trực quan hóa các kết quả nghiên cứu phức tạp (Chương 5).
\item
  \textbf{Xu hướng tương lai}: Những nhận định mang tính dự báo nhưng dựa trên bằng chứng về việc AI có thể tiếp tục định hình giới học thuật như thế nào trong những năm tới (Chương 6).
\item
  \textbf{Mẹo thực hành và tài nguyên}: Hướng dẫn tận dụng AI hiệu quả đồng thời duy trì sự nghiêm cẩn học thuật và tư duy phản biện (Chương 7).
\item
  \textbf{Thách thức về quy định và đạo đức}: Các cuộc tranh luận về pháp lý, đạo đức và thể chế xung quanh việc sử dụng AI trong công việc học thuật, bao gồm các vấn đề về quyền tác giả, quyền riêng tư dữ liệu và thông tin sai lệch (Chương 8).
\end{itemize}

Cuốn sách này được thiết kế cho đông đảo độc giả trong giới học thuật, bao gồm:

\begin{itemize}
\item
  \textbf{Nhà nghiên cứu} muốn tích hợp AI vào quy trình làm việc của mình, từ tổng hợp tài liệu đến phân tích dữ liệu.
\item
  \textbf{Nhà giáo dục} đang tìm kiếm những cách thức sáng tạo để nâng cao sự tham gia của người học và tối ưu hóa tài liệu giảng dạy.
\item
  \textbf{Nghiên cứu sinh và học viên sau đại học} đang đối mặt với thách thức kép: làm chủ nghiên cứu có AI hỗ trợ và đồng thời phát triển quan điểm phản biện của riêng mình.
\item
  \textbf{Nhà quản lý và hoạch định chính sách} cần hiểu những hàm ý mang tính thể chế của việc áp dụng AI trong giáo dục đại học.
\end{itemize}

Mặc dù hiểu biết về AI không phải là điều kiện tiên quyết để đọc cuốn sách này, chúng tôi vẫn cung cấp những giải thích nền tảng về các khái niệm then chốt nhằm bảo đảm tính dễ tiếp cận đối với mọi lĩnh vực. Mỗi chương của cuốn sách đi sâu vào những mảng cụ thể mà AI đang tạo ra tác động. Từ hỗ trợ hình thành ý tưởng nghiên cứu đến quản lý dữ liệu quy mô lớn, AI đang chứng tỏ là một cộng sự giá trị trong học thuật. Tuy nhiên, việc sử dụng AI cũng làm nảy sinh những mối quan ngại nghiêm trọng, đặc biệt liên quan đến tính minh bạch, sở hữu trí tuệ và nguy cơ làm suy giảm kỹ năng của nhà nghiên cứu.

Chúng tôi đưa ra một góc nhìn cân bằng về những lợi ích và hạn chế của AI bằng cách nêu bật các nghiên cứu tình huống thực tiễn và cung cấp các chiến lược thực hành nhằm tích hợp AI vào công việc học thuật. Nội dung này bao gồm:

\begin{itemize}
\item
  Cách AI có thể nâng cao, nhưng không thay thế, việc viết học thuật chặt chẽ.
\item
  Vai trò của AI trong xử lý và trực quan hóa dữ liệu, đặc biệt là khi xử lý các tập dữ liệu lớn, phi cấu trúc.
\item
  Các nguyên tắc đạo đức để sử dụng AI một cách có trách nhiệm, bảo đảm rằng AI bổ trợ thay vì làm tổn hại đến liêm chính học thuật.
\end{itemize}

\section{Các khái niệm nền tảng: Hiểu về AI tạo sinh}

Để tiếp cận AI tạo sinh (generative AI) một cách hiệu quả, điều quan trọng là phải hiểu các hệ thống này vận hành như thế nào, giới hạn của chúng là gì, và chúng khác với các loại trí tuệ nhân tạo (artificial intelligence) khác ra sao. AI trong môi trường học thuật thường được bàn đến theo nghĩa rộng, nhưng AI tạo sinh đặc biệt chỉ những mô hình có khả năng tạo ra văn bản, hình ảnh, mã nguồn, thậm chí cả âm nhạc giống như do con người tạo ra. Phần này phân tích các thành phần chính của AI tạo sinh, bao gồm cách nó được huấn luyện, các mô hình cốt lõi đang được sử dụng hiện nay, và những nền tảng lý thuyết thúc đẩy năng lực của nó.

AI tạo sinh là một nhánh con của trí tuệ nhân tạo, chuyên tạo ra nội dung mới thay vì chỉ phân loại, phân tích hoặc truy xuất thông tin hiện có. Không giống các mô hình AI truyền thống được thiết kế cho những tác vụ cụ thể, như phát hiện thư rác, nhận dạng hình ảnh hoặc hệ thống gợi ý, AI tạo sinh có thể sinh ra các đầu ra mới dựa trên những mẫu hình đã học. Nó sử dụng các phương pháp xác suất, dự đoán chuỗi từ, điểm ảnh hoặc sóng âm có khả năng xảy ra cao nhất để tạo ra nội dung mạch lạc và phù hợp ngữ cảnh.

AI tạo sinh hiện đại chủ yếu được thúc đẩy bởi học sâu (deep learning), một nhánh của học máy (machine learning) sử dụng các mạng nơ-ron nhân tạo nhiều tầng. Các mạng này xử lý lượng dữ liệu khổng lồ để nhận diện các mối quan hệ, cấu trúc và sắc thái ngữ cảnh, giúp chúng bắt chước nội dung do con người tạo ra. Tuy nhiên, cần nhấn mạnh rằng AI không ``hiểu'' như con người. Thay vào đó, nó vận hành dựa trên các tương quan thống kê, ánh xạ đầu vào thành đầu ra có khả năng cao nhất mà không sở hữu tri thức nội tại, lý luận hay chủ đích.

\subsection{Các mô hình và kiến trúc chính: ``Bộ não'' đằng sau AI tạo sinh}

Một số mô hình AI lớn đã nổi lên như những lực lượng chủ đạo trong AI tạo sinh, mỗi mô hình có kiến trúc và ứng dụng riêng biệt. Việc hiểu các mô hình này giúp làm sáng tỏ cách nội dung do AI tạo ra được hình thành, cũng như điểm mạnh và hạn chế của chúng nằm ở đâu.

\subsubsection{1. Các mô hình dựa trên Transformer: Xương sống của AI hiện đại}

Bước đột phá giúp AI tạo sinh ngày nay ra đời là sự xuất hiện của \textbf{kiến trúc transformer}. Không giống các mô hình mạng nơ-ron trước đây xử lý đầu vào theo trình tự (chẳng hạn mạng nơ-ron hồi quy, RNN), transformer sử dụng \textbf{cơ chế tự chú ý (self-attention)} để phân tích song song toàn bộ chuỗi đầu vào. Điều này cho phép mô hình nhận biết các phụ thuộc tầm xa trong văn bản, qua đó cải thiện đáng kể khả năng tạo ra các phản hồi mạch lạc, nhạy bén với ngữ cảnh.

Các mô hình dựa trên transformer được sử dụng rộng rãi nhất gồm:

\begin{itemize}
\item
  \textbf{GPT (Generative Pre-trained Transformer)} - Được OpenAI phát triển, họ mô hình này (GPT-3, GPT-4 và các phiên bản tương lai) được huấn luyện trên các kho ngữ liệu văn bản khổng lồ để tạo ra văn bản trôi chảy, giống con người. Các mô hình GPT được dùng rộng rãi trong viết học thuật, lập trình và AI hội thoại.
\item
  \textbf{Claude} - Do Anthropic tạo ra, Claude tập trung vào tính an toàn, khả năng diễn giải và giảm thiểu thiên lệch, khiến nó trở thành một lựa chọn thay thế mạnh mẽ cho việc sử dụng AI có trách nhiệm trong môi trường học thuật.
\item
  \textbf{BERT (Bidirectional Encoder Representations from Transformers)} - Không giống GPT, vốn thiên về tạo sinh, BERT chủ yếu được thiết kế để hiểu và phân loại văn bản. Mô hình này vượt trội ở các tác vụ như trả lời câu hỏi và tóm tắt văn bản.
\item
  \textbf{Gemini (bộ công cụ AI của Google, trước đây là Bard)} - Một AI đa phương thức tích hợp dữ liệu văn bản, hình ảnh và âm thanh nhằm nâng cao khả năng thích ứng của AI tạo sinh.
\item
  \textbf{Llama (Large Language Model Meta AI của Meta)} - Một lựa chọn mã nguồn mở, hiệu năng cao thay thế cho các mô hình độc quyền, được thiết kế để dễ tùy biến và mở rộng quy mô.
\item
  \textbf{DeepSeek} - Một lựa chọn mã nguồn mở của Trung Quốc, có hiệu quả cao, thay thế cho các mô hình độc quyền, sử dụng cái gọi là ``chú ý tiềm ẩn đa đầu (multi-head latent attention)'' để giảm kích thước bộ nhớ đệm khóa-giá trị, cho phép các truy vấn tiêu tốn ít tài nguyên hơn nhiều.
\end{itemize}

Tất cả các mô hình này đều có chung cấu trúc dựa trên transformer, nhưng mục tiêu huấn luyện của chúng khác nhau. Các mô hình GPT tạo văn bản theo kiểu tự hồi quy (dự đoán từ tiếp theo dựa trên ngữ cảnh trước đó). Ngược lại, BERT được tối ưu hóa cho việc hiểu hai chiều, khiến nó phù hợp hơn với các tác vụ hiểu ngôn ngữ thay vì tạo sinh văn bản.

\subsubsection{2. Huấn luyện AI tạo sinh: Từ dữ liệu thô đến đầu ra thông minh}

Các mô hình AI tạo sinh được huấn luyện qua ba giai đoạn chính:

\begin{enumerate}
\def\labelenumi{\arabic{enumi}.}
\item
  \textbf{Tiền huấn luyện (Pre-training)} - Các mô hình được tiếp xúc với những tập dữ liệu khổng lồ (ví dụ: sách, bài báo học thuật, Wikipedia và các diễn đàn internet) để học các mẫu ngôn ngữ chung, cú pháp và các liên hệ về mặt thực tế. Trong giai đoạn này, chúng dự đoán những từ bị thiếu trong câu, qua đó cải thiện khả năng tạo ra ngôn ngữ mạch lạc.
\item
  \textbf{Tinh chỉnh (Fine-tuning)} - Sau khi tiền huấn luyện, các mô hình được hoàn thiện thêm bằng những tập dữ liệu có mục tiêu, nhằm bảo đảm phù hợp với các trường hợp sử dụng cụ thể. Chẳng hạn, một mô hình AI được tinh chỉnh cho việc viết học thuật có thể được huấn luyện trên các bài báo tạp chí và kỷ yếu hội nghị để nâng cao độ chính xác theo lĩnh vực chuyên môn.
\item
  \textbf{Học tăng cường từ phản hồi của con người (Reinforcement Learning from Human Feedback, RLHF)} - Quy trình này nâng cao hiệu suất của AI bằng cách đưa các đánh giá của con người về những phản hồi được tạo ra vào quá trình huấn luyện. Những người huấn luyện xếp hạng các đầu ra do AI tạo ra, củng cố các hành vi mong muốn như độ chính xác về mặt thực tế, tính mạch lạc và các cân nhắc đạo đức.
\end{enumerate}

Ngay cả với những bước huấn luyện này, các mô hình AI tạo sinh vẫn mang tính \textbf{ngẫu nhiên (stochastic)}, đầu ra của chúng mang tính xác suất chứ không mang tính tất định. Điều này có nghĩa là cùng một prompt (câu lệnh) có thể cho ra các phản hồi khác nhau mỗi lần, khiến nội dung do AI tạo ra hữu ích cho việc hình thành ý tưởng nhưng thiếu tin cậy về tính nhất quán thực tế nếu không có sự giám sát của con người.

\subsection{Cách AI ``suy nghĩ'': Cơ chế tạo sinh ngôn ngữ}

Một quan niệm sai lầm phổ biến cho rằng AI tạo văn bản bằng cách truy xuất các câu trả lời đã được viết sẵn. Trên thực tế, các LLM không ``cắt/sao chép/dán'' từ dữ liệu huấn luyện của chúng, mà tạo văn bản một cách linh hoạt dựa trên các phân phối xác suất. Quy trình này diễn ra như sau:

\begin{enumerate}
\def\labelenumi{\arabic{enumi}.}
\item
  \textbf{Tách token (Tokenization)} - Văn bản được chia thành các đơn vị nhỏ hơn gọi là \textbf{token}, có thể là từ, từ con hoặc thậm chí là ký tự. Chẳng hạn, cụm từ ``machine learning'' (học máy) có thể được tách thành {[}``machine'', ``learning''{]} hoặc {[}``mach'', ``ine'', ``learning''{]} tùy thuộc vào bộ tách token được sử dụng.
\item
  \textbf{Phân tích ngữ cảnh (Contextual Analysis)} - Bằng cơ chế tự chú ý (self-attention), mô hình xác định những từ nào trong một chuỗi có liên quan nhất để dự đoán token tiếp theo.
\item
  \textbf{Dự đoán token tiếp theo (Next-Token Prediction)} - AI chọn từ tiếp theo có xác suất thống kê cao nhất dựa trên dữ liệu huấn luyện và đầu vào được cung cấp. Nếu có sự mơ hồ, nó tạo ra đầu ra dựa trên \textbf{thiết lập nhiệt độ (temperature)} (một tham số kiểm soát mức độ ngẫu nhiên, giá trị thấp cho ra câu trả lời mang tính xác định, trong khi giá trị cao khuyến khích sự biến thể sáng tạo).
\item
  \textbf{Lặp lại và hoàn tất (Iteration and Completion)} - Quy trình lặp lại cho đến khi mô hình đạt đến điểm dừng được xác định trước (chẳng hạn như giới hạn về số câu hoặc đoạn văn).
\end{enumerate}

Phương pháp này giải thích vì sao văn bản do AI tạo ra đôi khi sai một cách quá tự tin, nó không kiểm chứng sự kiện mà chỉ xây dựng các chuỗi nghe có vẻ hợp lý dựa trên các khuôn mẫu trước đó.

\section{Bối cảnh lịch sử: Từ AI biểu tượng đến cuộc cách mạng tạo sinh}

Lĩnh vực trí tuệ nhân tạo đã trải qua những biến động mạnh mẽ trong vài thập kỷ qua, đi qua các chu kỳ lạc quan cao độ, bùng nổ nguồn vốn, thất bại và hồi sinh. Sự trỗi dậy của AI tạo sinh là bước chuyển mình mới nhất trong quá trình tiến hóa liên tục này, nhưng nền tảng của nó đã được đặt ra bởi những đột phá trước đó trong lập luận biểu tượng, học máy và mạng nơ-ron. Hiểu lịch sử này là điều cốt yếu để đánh giá đúng cả thời điểm hiện tại lẫn quỹ đạo mà AI nhiều khả năng sẽ đi theo trong giới học thuật và xa hơn nữa.

Nguồn gốc của trí tuệ nhân tạo có thể truy về những năm 1950, khi các nhà nghiên cứu tìm cách tạo ra những cỗ máy có khả năng tái tạo lập luận của con người. Bài báo mang tính khai mở của Alan Turing, \emph{``Computing Machinery and Intelligence''} (1950), đã đặt ra câu hỏi \emph{``Máy móc có thể suy nghĩ không?''} và đưa ra Phép thử Turing nổi tiếng ngày nay như một thước đo cho trí tuệ nhân tạo. Cùng thời gian đó, John McCarthy, Marvin Minsky, Herbert Simon và Allen Newell đã phát triển các chương trình AI đầu tiên thực hiện lập luận biểu tượng, giải các câu đố logic và chứng minh các định lý toán học.

\begin{itemize}
\item
  Logic Theorist (1956) và General Problem Solver (1957) của Newell và Simon nhằm tự động hóa lập luận logic.
\item
  McCarthy đã đặt ra thuật ngữ Trí tuệ nhân tạo vào năm 1956 tại Hội nghị Dartmouth, sự kiện sáng lập của lĩnh vực này.
\item
  Các hệ thống xử lý ngôn ngữ tự nhiên (NLP) ban đầu, chẳng hạn ELIZA (1966) của Joseph Weizenbaum, cố gắng mô phỏng hội thoại của con người.
\end{itemize}

Trong thời kỳ này, nghiên cứu AI bị chi phối bởi các phương pháp biểu tượng, vốn dựa trên những quy tắc được lập trình tường minh để giải quyết vấn đề. Tuy nhiên, các hệ thống ban đầu này gặp khó khăn trước sự phức tạp của thế giới thực: chúng hoạt động tốt trong những môi trường bị giới hạn (như chơi cờ vua) nhưng thất bại ở các nhiệm vụ mở đòi hỏi lập luận theo lẽ thường.

``Mùa đông AI'' đầu tiên (1973--những năm 1980) được châm ngòi bởi sự hoài nghi ngày càng tăng đối với những lời hứa lớn lao và các hạn chế thực tiễn của AI. Hai báo cáo quan trọng đã thổi bùng sự hoài nghi này:

\begin{enumerate}
\def\labelenumi{\arabic{enumi}.}
\item
  Báo cáo Lighthill (1973) ở Anh lập luận rằng AI đã không đạt được các mục tiêu đầy tham vọng của mình và khó có thể biện minh cho việc tiếp tục được chính phủ tài trợ.
\item
  DARPA (Cơ quan Dự án Nghiên cứu Quốc phòng Tiên tiến Hoa Kỳ) cũng cắt giảm tài trợ, kết luận rằng nghiên cứu AI đang cho lợi tức giảm dần.
\end{enumerate}

Kết quả là nhiều dự án AI mất nguồn tài trợ và nghiên cứu chậm lại. Những thách thức chính là:

\begin{itemize}
\item
  Hạn chế về tính toán: Các mô hình AI ban đầu đòi hỏi sức mạnh xử lý khổng lồ mà thời đó chưa có.
\item
  Thiếu khả năng thích ứng: Các hệ thống AI không thể khái quát hóa tri thức ra ngoài những quy tắc hẹp, định sẵn.
\item
  Thất bại của AI biểu tượng: Các hệ thống dựa trên quy tắc không thể xử lý dữ liệu mơ hồ, mang tính xác suất hoặc không đầy đủ.
\end{itemize}

Bất chấp suy thoái này, những hạt giống của sự hồi sinh AI trong tương lai vẫn lặng lẽ được gieo. Cuối thập niên 1980, các nhà nghiên cứu bắt đầu chuyển trọng tâm từ AI biểu tượng sang các mô hình kết nối (connectionist), về sau trở thành nền tảng của học sâu.

Đến thập niên 1990, những tiến bộ về sức mạnh tính toán, mô hình hóa thống kê và khả năng sẵn có của dữ liệu đã dẫn đến sự phục hưng của nghiên cứu AI, đánh dấu bước chuyển từ AI dựa trên quy tắc sang các phương pháp dựa trên dữ liệu.

\begin{itemize}
\item
  Học máy nổi lên như mô hình thống trị, chuyển các hệ thống AI từ những quy tắc được lập trình tường minh sang học từ dữ liệu.
\item
  Máy vector hỗ trợ (Support Vector Machines, SVM) và Mô hình Markov ẩn (Hidden Markov Models, HMM) tạo ra những đột phá trong nhận dạng mẫu và xử lý giọng nói.
\item
  Mạng nơ-ron được hồi sinh nhờ các thuật toán huấn luyện được cải thiện, chẳng hạn lan truyền ngược (backpropagation) (Hinton, 1986).
\end{itemize}

Bước chuyển này cho phép các ứng dụng thực tiễn như lọc thư rác, hệ thống gợi ý và các trợ lý giọng nói ban đầu (ví dụ: Watson của IBM). AI bắt đầu rời phòng thí nghiệm nghiên cứu để đi vào các ứng dụng thực tế.

Thập niên 2010 chứng kiến sự bùng nổ về năng lực của AI, được thúc đẩy bởi ba yếu tố chính:

\begin{enumerate}
\def\labelenumi{\arabic{enumi}.}
\item
  Sức mạnh tính toán chưa từng có - Sự ra đời của Bộ xử lý đồ họa (GPU), vốn được thiết kế cho trò chơi điện tử, đã cho phép huấn luyện mạng nơ-ron quy mô lớn.
\item
  Sự sẵn có của Dữ liệu lớn - Internet, mạng xã hội và các nguồn học thuật được số hóa đã cung cấp cho AI những tập dữ liệu khổng lồ để học hỏi.
\item
  Các thuật toán đột phá - Việc phát triển các mạng nơ-ron sâu, đặc biệt là Mạng nơ-ron tích chập (CNN) cho nhận dạng hình ảnh và Mạng nơ-ron hồi quy (RNN) cho mô hình hóa ngôn ngữ, đã làm thay đổi lĩnh vực này.
\end{enumerate}

Một số cột mốc quan trọng gồm:

\begin{itemize}
\item
  2012 - Mô hình AlexNet giành chiến thắng trong cuộc thi ImageNet, chứng minh sức mạnh của học sâu đối với nhận dạng thị giác.
\item
  2014 - Sự ra đời của Mạng đối sinh (Generative Adversarial Networks, GAN) (Goodfellow et al.), cho phép AI tạo ra hình ảnh chân thực, deepfake và nội dung nghệ thuật.
\item
  2017 - Sự xuất hiện của các mô hình transformer (\emph{``Attention Is All You Need''}, Vaswani et al.), mở đường cho các mô hình ngôn ngữ lớn như GPT.
\end{itemize}

Với sự xuất hiện của các kiến trúc dựa trên transformer, AI tạo sinh bước vào một giai đoạn tinh vi mới. Không giống các mô hình trước đây vốn phân loại hoặc dự đoán dựa trên những khuôn mẫu có sẵn, các hệ thống mới này có thể tạo ra văn bản, mã nguồn, hình ảnh mạch lạc, nguyên bản, thậm chí cả giọng nói tổng hợp.

Các bước phát triển then chốt của LLM (Mô hình ngôn ngữ lớn) gồm:

\begin{itemize}
\item
  GPT-3 (2020) - Một cột mốc trong văn bản do AI tạo ra, có khả năng đưa ra những phản hồi tinh vi và viết sáng tạo.
\item
  DALL-E (2021) - Một mô hình tạo ra hình ảnh do AI sáng tạo từ các prompt (câu lệnh) dạng văn bản, đánh dấu một bước nhảy vọt của AI đa phương thức.
\item
  ChatGPT (2022) - Đưa AI tạo sinh vào sử dụng đại chúng, cách mạng hóa tương tác giữa con người và AI.
\item
  Claude, Llama và Gemini (2023--2024) - Các mô hình AI cạnh tranh mang lại những tiến bộ về lập luận, tính mạch lạc theo ngữ cảnh và các biện pháp bảo vệ đạo đức.
\item
  {[}DeepSeek{[} (2025) - Một mô hình mang lại những cải thiện đáng kể về hiệu quả, đặt câu hỏi về nhu cầu sử dụng phần cứng đắt tiền để đạt được kết quả hàng đầu.{]}{]}
\end{itemize}

Lịch sử AI cho thấy những chu kỳ lặp đi lặp lại của sự cường điệu và vỡ mộng. Dù năng lực của AI ngày nay rõ ràng mang tính chuyển đổi, vẫn tồn tại những rủi ro dai dẳng:

\begin{itemize}
\item
  {[}Hứa hẹn quá mức so với thực tế - Các Mùa đông AI trong quá khứ dạy chúng ta rằng những kỳ vọng thổi phồng có thể dẫn đến phản ứng ngược khi AI không thực hiện được những tuyên bố quá tham vọng.{]}
\item
  {[}Hợp tác giữa con người và AI - AI không thay thế các nhà nghiên cứu; nó nâng cao khả năng làm việc hiệu quả, tư duy phản biện và tiếp cận những ý tưởng phức tạp của họ.{]}
\item
  {[}Nhận thức về đạo đức và quy định - Sự phát triển nhanh chóng của AI đòi hỏi các cuộc thảo luận liên tục về việc sử dụng có trách nhiệm, tính minh bạch và quản trị.{]}
\end{itemize}

\section{Con đường phía trước}

Khi tiến về phía trước, AI tạo sinh sẽ tiếp tục định hình sâu sắc công việc học thuật. Thách thức đối với các học giả, nhà giáo dục và các tổ chức không phải là có nên sử dụng AI hay không, mà là sử dụng AI như thế nào một cách có trách nhiệm. Cuốn sách này cung cấp một sự khám phá có nền tảng lịch sử và mang tính phản biện về tiềm năng của AI trong môi trường học thuật, bảo đảm rằng việc tích hợp AI phù hợp với các giá trị học thuật, các cân nhắc đạo đức và sự nghiêm cẩn trí tuệ.

Thay vì xem AI như một công cụ trung lập, chúng tôi khuyến khích các nhà nghiên cứu tiếp cận những công nghệ này một cách có tính phản biện. AI không chỉ đơn thuần là một công cụ tăng hiệu suất, mà còn là một lực lượng định hình bản chất của việc sản xuất và phổ biến tri thức. Bằng cách hiểu cơ chế vận hành của AI, tận dụng điểm mạnh và giảm thiểu điểm yếu của nó, các học giả có thể khai thác tiềm năng của AI đồng thời duy trì tính toàn vẹn của hoạt động nghiên cứu học thuật.

Cuốn sách này không đưa ra một công thức chung cho mọi trường hợp về việc sử dụng AI trong học thuật. Thay vào đó, sách cung cấp một nền tảng cho việc ra quyết định có hiểu biết, trang bị cho độc giả kiến thức và chiến lược cần thiết để định hướng trong lĩnh vực đang thay đổi nhanh chóng này. Khi chúng ta bắt đầu hành trình tìm hiểu vai trò của AI trong học thuật, chúng tôi mời độc giả tiếp cận với sự tò mò, sự thận trọng và cam kết đối với nghiên cứu học thuật có đạo đức.
